% Pursuing Stacks (cotes 134-2 to 134-8): archive cote 134-5, site page 63 (PDF page 64), typescript page 250.
% typescript. Draft compared: ps4.tex:3172-3233.
% AI-transcribed by Claude Opus 5.5 (claude-opus-5-5, adaptive thinking, effort high), Message Batches API, batch ps2 (msgbatch_01NVwr1Yuo8qJf5rykeGBXPK), 2026-11-12; reviewed against both supplied scan halves; see CHECK_DECISIONS.tsv.
\scanpage{134-5/63}
\mspage{250}

I better correct the sta\inserted{t}ement, to the effect that \inserted{it seems} all modelizing \hand{\uline{\textcolor{black}{totally}}}
\hand{\uline{\textcolor{black}{aspheric}}} asphericity structures met with so far \deleted{seem to be} \inserted{are} generated by
suitable contractibility structures.

Let M be a category endowed with a contractibility structure $M_c$ ,
hence also with an asphericity structure $M_{as}$ , and let \xout{\unsure{ixxAxxxxMxxbexxx}}
\xout{\unsure{xxxx}}
\xout{functor from a small category \unsure{IxxAxxxxMx}} \xout{\unsure{Ixx}}
\[ i : A \hand{\to} M \]
be a functor from a small category A to M. We want to give ns conditions
\over{f}{f}or i to be a \hand{\uline{\textcolor{black}{test}}} \hand{\uline{\textcolor{black}{functor}}}, in terms of homotopy notions in M and in A\^{}.
To this end, it seems necessary to refresh \hand{(}somewhat memory\hand{)}
% A continuous handwritten transposition stroke runs beneath somewhat and above memory; the typed word order is preserved.
 and recall
some relevant notions which were developped in part IV of our notes(sec-%
\hand{\uline{\textcolor{black}{tions 54 and 55).}}}

It will be convenient to call an object \inserted[typed]{F} of A\^{} \hand{\uline{\textcolor{black}{locally}}} \hand{\uline{\textcolor{black}{aspheric}}}
(resp. \hand{\uline{\textcolor{black}{totally}}} \hand{\uline{\textcolor{black}{aspheric}}} ), if it's product \inserted[typed]{in A\^{}} by any object of A, and hence
also it's product by any aspheric object of A\^{}, is aspheric (resp. and
if moreover F itself is aspheric). With this terminology, A\^{} is totally
aspheric iff every aspheric object of A\^{} is totally aspheric, and \inserted[typed]{if} more-%
over A is aspheric i.e. $e_{A\hat{}}$ is aspheric. \xout{If A is aspheric\unsure{xxxthenxIx}} Note
that \xout{\unsure{the}} \inserted[typed]{F is} locally aspheric \xout{\unsure{objectxx}} iff the map $F \hand{\to} e_{A\hat{}}$ is aspheric ,
or what amounts to the same, if this map is universally in $\underline{W}_A$ , or
equivalently, if the corresponding functor $A_{/F} \hand{\to} A$ is aspheric.
If A itself is aspheric, and\hand{/}in this case only, this condition implies
already that F is aspheric, i.e. that F is totally aspheric. We'll
denote by \xout{\unsure{it}}
\[ \hat{A}_{tot.as} \quad ( \text{ resp. } \quad \hat{A}_{loc.\ as} \ ) \]
the full subcategory of A\^{} made up with the totally aspheric (resp.
locally aspheric) objects of A\^{} . Thus we get
% Hand-drawn arrows are marked blue; typed objects remain plain.
\diagram{
\begin{tikzcd}[column sep=small]
 & \hat{A}_{tot.as} = \hat{A}_{loc.as} \hand{\cap} \hat{A}_{as} \ar[dl, hook', color=blue] \ar[dr, hook, color=blue] & \\
\hat{A}_{loc.as} & & \hat{A}_{as} \ .
\end{tikzcd}
}

N\over{o}{\unsure{o}}w recall that (section 54) in terms of the set $\underline{W}_A$ of weak equi-%
valences in A\^{}, we constructed a homotopy structure on A\^{}, more speci-%
fically a homotopy interval structure, admitting as a generating family
of homotopy intervals \xout{\unsure{f}}
\[ \underline{I} =(I,\hand{\delta}_0,\hand{\delta}_1) \]
the set of intervals such that I be a locally aspheric object of A\^{}\hand{\unsure{,}}\xout{\unsure{xi}}i.e.
% The uncertain comma-like tick is raised above the type line near i.e.; its function is unresolved.
\[ I \hand{\in} \mathrm{Ob}\, \hat{A}_{loc.\ as} \quad . \]

\margin{63}
